\documentclass[12pt]{article}
\usepackage[utf8]{inputenc}
\usepackage[T1]{fontenc}
\usepackage{hyperref}
\title{Kagoro Language Resources for K-12 Curriculum Development}
\author{Compiled by Claude Sonnet 5.5 (claude-sonnet-5-5) for Andr\'as Kornai}
\date{}
\begin{document}
\maketitle

\section{Introduction}
Kagoro (also Kakolo; ISO 639-3 xkg; Glottocode kago1248) is a Western Mande language of the Manding group, spoken in central Mali (Kaarta-Bin\'e and Gumbu areas). The speaker count is uncertain. About 15,000 is the figure in the brief. EndangeredLanguages.com cites 21,500 from a 1998 source. A 2001 SIL survey found mixed vitality, and Bambara is gradually replacing Kagoro. The search results show lexical similarity of about 86\% with Kita Maninka. In the search results, nothing shows Kagoro being taught in schools. The material found is thin and consists only of general descriptive pages.

\section{Primary Educational Resources}
No primers, textbooks or school programmes for Kagoro were found in the search results. A curriculum would have to be built from scratch. It could draw on Bambara and Maninka materials, which are widely used in Mali, but no such materials were identified in this search.

\section{Linguistic and Reference Materials}
\begin{itemize}
\item Endangered Languages Project entry for Kagoro: \url{https://endangeredlanguages.com/lang/3250?hl=fr}
\item Verbix endangered-languages page for Kagoro: \url{https://docs.verbix.com/EndangeredLanguages/Kagoro}
\item Wikipedia, Kita Maninka language (a closely related language that shares the xkg code in the ISO entry): \url{https://en.wikipedia.org/wiki/Kita_Maninka_language}
\end{itemize}
No dictionary, grammar or orthography for Kagoro was found.

\section{Religious and Community Resources}
\begin{itemize}
\item Joshua Project profile, Kagoro in Mali (demographic and religious overview only): \url{https://joshuaproject.net/people_groups/12413/ML}
\end{itemize}
No scripture translations were found in the search results.

\section{Digital and Audio Resources}
No audio, video or corpus resources for Kagoro were found.

\section{Curriculum Development Context}
Kagoro appears to be a vulnerable language under pressure from Bambara. Any curriculum work would need fieldwork first: community consultation, orthography design, and documentation. The related Manding languages (Bambara, Maninka) could provide a model for materials. The available web evidence is too thin to support more than this.
\end{document}